\documentclass[10pt,a4paper]{amsart}
\usepackage[margin=19mm]{geometry}
\usepackage{amsmath,amssymb,mathtools}
\usepackage[scr=rsfs,cal=euler]{mathalfa}
\usepackage{xfrac}
\usepackage[unicode,hidelinks,bookmarks=false]{hyperref}
\newcommand{\ie}{i.e.\ }
\newcommand{\cf}{cf.\ }
\newcommand{\resp}{respectively\ }
\newcommand{\Subsection}{\subsection}
\providecommand{\nobreakdash}{\nobreak\hbox{-}}
\setlength{\emergencystretch}{2em}
\multlinegap0pt
\usepackage{bm}
\usepackage{bbm}
\usepackage{dsfont}
\usepackage{xspace}
\usepackage{cancel}
\usepackage{mathtools}
\usepackage{amsthm}
\makeatletter
\@ifundefined{theorem}{\newtheorem{theorem}{Theorem}}{}
\@ifundefined{lemma}{\newtheorem{lemma}[theorem]{Lemma}}{}
\makeatother
\title[Traveling waves in nonclassical diffusion equations]{Traveling waves in nonclassical diffusion equations}
\author{William Barker, Le Xuan Dong, Vu Trong Luong, Nguyen Duong Toan}
\date{Source: arXiv:2510.22349v1 (CC BY 4.0)}
\begin{document}
\maketitle

\noindent\emph{Source and licence.} Traveling waves in nonclassical diffusion equations. Version arXiv:2510.22349v1 (CC BY 4.0), distributed under the Creative Commons Attribution 4.0 International licence, as recorded in the arXiv licence metadata for this exact version and shown on the abstract page https://arxiv.org/abs/2510.22349v1. Full rights notice: licenses/arxiv-2510.22349-CC-BY-4.0.txt.\par\smallskip

\noindent\textit{Editorial scope.} Counted unit: the Theorem whose source label is \texttt{smooth1} in the article (source0.tex lines 1551--1564, source label \texttt{smooth1}), delivered together with the article's own complete original proof of it (source0.tex lines 1565--1655, one \texttt{proof} environment of 91 lines). No mathematics is written, repaired or re-derived: statement and proof are the article's own text line for line. Required context retained uncounted: Green (lines 850--870). Every definition, display and label of those lines is retained unchanged. Omitted from this package and not redistributed: 1--849, 871--1550, 1656--2168. Nothing inside a retained excerpt is omitted or altered: every retained body is delivered exactly as the article's lines are written, so no omission receipt of any kind accompanies this package. Citations: the retained excerpts carry no citation command at all, so no bibliography entry of the article is transcribed and no reference is redistributed. Reference adaptation: none was needed and none was made; every reference inside the delivered unit resolves inside it, so no unresolved reference or citation is produced. Printed slips kept literal: no printed word, symbol, hypothesis or number of the counted statement or of its proof was altered. Layout and class: the article's own class is \texttt{amsart} with packages including hyperref, amssymb, xcolor, geometry; the wrapper is the lane's pinned \texttt{amsart} editorial wrapper at 10pt on A4, which restates the theorem-like environments the retained text needs and adds the article's own macro definitions that the retained text needs (none), with extra wrapper packages (bm, bbm, dsfont, xspace, cancel, mathtools). The delivered numbering is the unit's own. Assets: no retained span contains a figure environment, a graphics-inclusion command, a PDF-page inclusion, a table or a TikZ picture; no image file is copied into this package, the package declares no asset dependency and no image file is redistributed with it. Licence: version arXiv:2510.22349v1 of this article is distributed under the Creative Commons Attribution 4.0 International licence, as recorded in the arXiv licence metadata for this exact version and shown on the abstract page https://arxiv.org/abs/2510.22349v1. A full rights notice is delivered at licenses/arxiv-2510.22349-CC-BY-4.0.txt. Characters: every retained excerpt is pure ASCII, so the delivered engine, XeLaTeX with its default Latin Modern text font, reports no missing character.
  \begin{lemma}\label{Green}
  	Let $G(\xi)$ be the Green's function of $\mathcal L$ satisfying the jump/continuity
  	conditions \textup{(M1)--(M3)}.
  	Assume that the characteristic polynomial
  	\[
  	\Delta(\lambda)=c\lambda^3 + D\lambda^2 - c\lambda - \beta
  	\]
  	has three distinct real roots $\lambda_1<\lambda_2<0<\lambda_3$ and that $c>0$.
  	Then $G$ admits the representation
  	\[
  	G(\xi)=
  	\begin{cases}
  		A_3 e^{\lambda_3\xi}, & \xi<0,\\[4pt]
  		- \bigl(A_1 e^{\lambda_1\xi}+A_2 e^{\lambda_2\xi}\bigr), & \xi>0,
  	\end{cases}
  	\]
  	where $A_1,A_2,A_3$ are determined by (M1)--(M3). Moreover,
  	\[
  	G(\xi)<0 \quad \text{for all }\xi\in\mathbb{R}.
  	\]
  \end{lemma}

 \begin{theorem}\label{smooth1}
 	Let $\Gamma$ and $F:=-\mathcal L^{-1}H$ be as above, and assume that the Green
 	function $G$ of $\mathcal L$ satisfies the exponential decay estimate
 	$|G^{(k)}(x)|\le C_k e^{-\gamma|x|}$ for $k=0,1,2$ and some constants $C_k,\gamma>0$.
 	Let $\overline\phi,\underline\phi\in\Gamma$ be a supersolution and a subsolution,
 	respectively. Then the following regularity assertions hold and the ordering is preserved:
 	\begin{enumerate}
 		\item $F(\overline\phi),\,F(\underline\phi)\in \Gamma\cap W^{2,\infty}(\mathbb R)$.
 		\item $F^2(\overline\phi),\,F^2(\underline\phi)\in \Gamma\cap W^{3,\infty}(\mathbb R)$.
 		\item $F^3(\overline\phi),\,F^3(\underline\phi)\in \Gamma\cap BC^3(\mathbb R)$.
 	\end{enumerate}
 	In each case the inequalities (supersolution $\ge F(\cdot)\ge$ subsolution, etc.)
 	remain valid.
 \end{theorem}

 \begin{proof}
 	We first recall the standing hypotheses used below:
 	\begin{itemize}
 		\item $\mathcal L\phi=c\phi''' + D\phi'' - c\phi' -\beta\phi$ is the linear operator introduced in (5.5), with $\beta>0$.
 		\item $G$ denotes the Green kernel of $\mathcal L$, so that $\mathcal L G=\delta$, and we assume the exponential decay / integrability property
 		$$G,G',G''\in L^1(\mathbb R),\qquad |G^{(k)}(x)|\le C_k e^{-\gamma|x|},\ k=0,1,2.$$
 		\item $H:\Gamma\to BC(\mathbb R)$ is the nonlinear operator $H(\phi)(t)=f_c(\phi_t)+\beta\phi(t)$; by (C2),(C3) we have $H(\phi)\in L^\infty(\mathbb R)$ for every $\phi\in\Gamma$.
 		\item Lemma~\ref{Green} (proved earlier) yields $G(\xi)<0$ for all $\xi\in\mathbb R$, hence the convolution operator $-\mathcal L^{-1}$ is order-preserving on $BC(\mathbb R)$.
 	\end{itemize}
 	
 	Let $\overline\phi\in\Gamma$ be a supersolution (the argument for a subsolution is analogous). Define
 	\[
 	F(\overline\phi)(t) := -\mathcal L^{-1}H(\overline\phi)(t) \;=\; -\int_{\mathbb R} G(t-s)\,H(\overline\phi)(s)\,ds.
 	\]
 	We prove that $F(\overline\phi)\in W^{2,\infty}(\mathbb R)$ and $F(\overline\phi)\in\Gamma$.
 	
 %	\medskip\noindent\textbf{Step 1. Boundedness of the right-hand side.}
 	By hypothesis $\overline\phi\in\Gamma\subset BC(\mathbb R,[0,K])$ and by (C2),(C3)
 	the composite $H(\overline\phi)$ is bounded; set $M:=\|H(\overline\phi)\|_\infty<\infty$.
 	Since $G\in L^1(\mathbb R)$, the convolution defining $F(\overline\phi)$ is well defined
 	and satisfies
 	\[
 	\|F(\overline\phi)\|_\infty \le \|G\|_{L^1}\,\|H(\overline\phi)\|_\infty \le \|G\|_{L^1} M.
 	\]
 	
 	%\medskip\noindent\textbf{Step 2. Differentiability and $W^{2,\infty}$-bounds.}
 	Because $G',G''\in L^1(\mathbb R)$ and $H(\overline\phi)\in L^\infty(\mathbb R)$,
 	we may differentiate under the integral sign (dominated convergence). For almost every $t$,
 	\[
 	(F(\overline\phi))'(t) = -\int_{\mathbb R} G'(t-s)\,H(\overline\phi)(s)\,ds,
 	\]
 	\[
 	(F(\overline\phi))''(t) = -\int_{\mathbb R} G''(t-s)\,H(\overline\phi)(s)\,ds.
 	\]
 	Taking suprema and using Young's (convolution) estimate yields the uniform bounds
 	\[
 	\|(F(\overline\phi))'\|_\infty \le \|G'\|_{L^1}\,M,\qquad
 	\|(F(\overline\phi))''\|_\infty \le \|G''\|_{L^1}\,M.
 	\]
 	Hence $F(\overline\phi)\in W^{2,\infty}(\mathbb R)$, as required.
 	
 %	\medskip\noindent\textbf{Step 3. Limits at infinity (preservation of boundary values).}
 	Because $\overline\phi\in\Gamma$ is nondecreasing with
 	$\lim_{t\to-\infty}\overline\phi(t)=0$, $\lim_{t\to+\infty}\overline\phi(t)=K$,
 	the history map $t\mapsto\overline\phi_t$ satisfies $\overline\phi_t\to\hat 0$ as $t\to-\infty$
 	and $\overline\phi_t\to\hat K$ as $t\to+\infty$. By continuity of $f_c$ we deduce
 	\[
 	\lim_{t\to-\infty} H(\overline\phi)(t)=H_-=f_c(\hat 0)+\beta\cdot 0=0,
 	\qquad
 	\lim_{t\to+\infty} H(\overline\phi)(t)=H_+=f_c(\hat K)+\beta K=\beta K.
 	\]
 	Since $H(\overline\phi)\in L^\infty$ and $G\in L^1$, dominated convergence gives
 	\[
 	\lim_{t\to\pm\infty} F(\overline\phi)(t)
 	= -\Bigl(\lim_{s\to\pm\infty} H(\overline\phi)(s)\Bigr)\int_{\mathbb R} G(y)\,dy.
 	\]
 	It remains to evaluate $\int_{\mathbb R}G(y)\,dy$. Integrating the identity
 	$\mathcal L G=\delta$ over $\mathbb R$ and using that boundary terms arising from total
 	derivatives vanish because of the exponential decay of $G$ and its derivatives, we obtain
 	\[
 	\int_{\mathbb R} \mathcal L G(y)\,dy = -\beta\int_{\mathbb R} G(y)\,dy = \int_{\mathbb R}\delta(y)\,dy = 1,
 	\]
 	hence
 	\[
 	\int_{\mathbb R} G(y)\,dy = -\frac{1}{\beta}.
 	\]
 	Therefore
 	\[
 	\lim_{t\to-\infty} F(\overline\phi)(t) = -H_- \int G = -0\cdot\Bigl(-\frac{1}{\beta}\Bigr)=0,
 	\]
 	\[
 	\lim_{t\to+\infty} F(\overline\phi)(t) = -H_+ \int G
 	= -(\beta K)\Bigl(-\frac{1}{\beta}\Bigr) = K.
 	\]
 	Thus $F(\overline\phi)$ attains the correct boundary values and so lies in the same
 	range of asymptotic states as $\Gamma$.
 	
 %	\medskip\noindent\textbf{Step 4. Monotonicity (order preservation).}
 	By hypothesis (C2)/(5.4) the map $H$ is order-preserving on $\Gamma$ (i.e.\ \(\phi\le\psi\)
 	implies \(H(\phi)\le H(\psi)\)). Since $G(\xi)<0$ for all $\xi$, the linear map
 	$h\mapsto -\int G(\cdot-s)h(s)\,ds=-\mathcal L^{-1}h$ is order-preserving on $BC(\mathbb R)$.
 	Hence $F=-\mathcal L^{-1}\circ H$ is order-preserving, so $F(\overline\phi)$ is nondecreasing
 	whenever $\overline\phi$ is. Combining monotonicity with the limits from Step~3 we conclude
 	$F(\overline\phi)\in\Gamma\cap W^{2,\infty}(\mathbb R)$.
 	
 	\medskip
 	The same argument applies to any subsolution $\underline\phi$. This completes the proof of part (i).
 	Parts (ii)--(iii) follow by the standard bootstrap: once $F(\phi)\in W^{2,\infty}$ the history
 	map $t\mapsto (F\phi)_t$ is continuous in $C([-\tau,0])$, so $H(F\phi)$ is continuous and bounded;
 	differentiating one more time under the integral gives $F^2(\phi)\in W^{3,\infty}$, etc.
 \end{proof}

\end{document}
